\documentclass[10pt,a4paper]{amsart}
\usepackage[margin=19mm]{geometry}
\usepackage{amsmath,amssymb,mathtools}
\usepackage[scr=rsfs,cal=euler]{mathalfa}
\usepackage{xfrac}
\usepackage[unicode,hidelinks,bookmarks=false]{hyperref}
\newcommand{\ie}{i.e.\ }
\newcommand{\cf}{cf.\ }
\newcommand{\resp}{respectively\ }
\newcommand{\Subsection}{\subsection}
\providecommand{\nobreakdash}{\nobreak\hbox{-}}
\setlength{\emergencystretch}{2em}
\multlinegap0pt
\usepackage{bm}
\usepackage{bbm}
\usepackage{dsfont}
\usepackage{xspace}
\usepackage{cancel}
\usepackage{mathtools}
\usepackage[shortlabels]{enumitem}
\usepackage{amsthm}
\makeatletter
\@ifundefined{theorem}{\newtheorem{theorem}{Theorem}}{}
\@ifundefined{lemma}{\newtheorem{lemma}[theorem]{Lemma}}{}
\makeatother
\def\norm#1#2{\left\|#2\right\|_{#1}} 
\def\snorm#1#2{\left|#2\right|_{#1}} 
\title[Error estimates of fully semi-Lagrangian schemes for diffusi]{Error estimates of fully semi-Lagrangian schemes for diffusive conservation laws}
\author{Haruki Takemura}
\date{Source: arXiv:2508.03455v2 (CC BY 4.0)}
\begin{document}
\maketitle

\noindent\emph{Source and licence.} Error estimates of fully semi-Lagrangian schemes for diffusive conservation laws. Version arXiv:2508.03455v2 (CC BY 4.0), distributed under the Creative Commons Attribution 4.0 International licence, as recorded in the arXiv licence metadata for this exact version and shown on the abstract page https://arxiv.org/abs/2508.03455v2. Full rights notice: licenses/arxiv-2508.03455-CC-BY-4.0.txt.\par\smallskip

\noindent\textit{Editorial scope.} Counted unit: the Theorem whose source label is \texttt{thm:1} in the article (source0.tex lines 415--425, source label \texttt{thm:1}), delivered together with the article's own complete original proof of it (source0.tex lines 427--646, one \texttt{proof} environment of 220 lines). No mathematics is written, repaired or re-derived: statement and proof are the article's own text line for line. Required context retained uncounted: SL-V (lines 104--110); p:o (lines 153--167); p:s (lines 155--157); p:q (lines 160--162); lem:s (lines 196--202); def\_tau (lines 259--262); tau\_0 (lines 349--351); tau\_s (lines 353--355); b7 (lines 362--364); b7-2 (lines 366--368); lem:2 (lines 372--377); lem:s2 (lines 1128--1133); lem:a1 (lines 1329--1339); lem:c (lines 1414--1425); d4 (lines 1501--1503). Every definition, display and label of those lines is retained unchanged. Omitted from this package and not redistributed: 1--103, 111--152, 158--159, 163--195, 203--258, 263--348, 352--352, 356--361, 365--365, 369--371, 378--414, 426--426, 647--1127, 1134--1328, 1340--1413, 1426--1500, 1504--1611. Nothing inside a retained excerpt is omitted or altered: every retained body is delivered exactly as the article's lines are written, so no omission receipt of any kind accompanies this package. Citations: the retained excerpts carry no citation command at all, so no bibliography entry of the article is transcribed and no reference is redistributed. Reference adaptation: none was needed and none was made; every reference inside the delivered unit resolves inside it, so no unresolved reference or citation is produced. Printed slips kept literal: no printed word, symbol, hypothesis or number of the counted statement or of its proof was altered. Layout and class: the article's own class is \texttt{amsart} with packages including geometry, graphicx, amsmath, amssymb, amsthm, mathtools, bm, comment, float, ascmac, color, cite, enumerate, multirow; the wrapper is the lane's pinned \texttt{amsart} editorial wrapper at 10pt on A4, which restates the theorem-like environments the retained text needs and adds the article's own macro definitions that the retained text needs (\texttt{\textbackslash norm}, \texttt{\textbackslash snorm}), with extra wrapper packages (bm, bbm, dsfont, xspace, cancel, mathtools, enumitem). The delivered numbering is the unit's own. Assets: no retained span contains a figure environment, a graphics-inclusion command, a PDF-page inclusion, a table or a TikZ picture; no image file is copied into this package, the package declares no asset dependency and no image file is redistributed with it. Licence: version arXiv:2508.03455v2 of this article is distributed under the Creative Commons Attribution 4.0 International licence, as recorded in the arXiv licence metadata for this exact version and shown on the abstract page https://arxiv.org/abs/2508.03455v2. A full rights notice is delivered at licenses/arxiv-2508.03455-CC-BY-4.0.txt. Characters: every retained excerpt is pure ASCII, so the delivered engine, XeLaTeX with its default Latin Modern text font, reports no missing character.
\begin{align} 
  \label{SL-V} 
  \begin{cases} 
    u_h^0 = \mathcal{I}_h u_0, & \\ 
    u_h^{n} = \mathcal{I}_h \left[\mathcal{S}_{\Delta t}^{\mathrm{V}} (\mathcal{S}_{\Delta t}^{\mathrm{A}} u_h^{n - 1})\right] & (n = 1, \ldots, N_t), 
  \end{cases} 
\end{align} 

\begin{enumerate}[(P1)] 
  \item The operator $ \mathcal{I}_{h} $ satisfies the integral relation \label{p:o} 
  \begin{align} 
    \snorm{s, 2}{(I - \mathcal{I}_{h}) v - \mathcal{I}_{h} w}^2 & = \snorm{s, 2}{(I - \mathcal{I}_{h}) v}^2 + \snorm{s, 2}{\mathcal{I}_{h} w}^2 
  \end{align} 
  for any $ v, w \in H^s (\mathbb{T}) $, where we denote the identity operator by $ I $. 
  \item There exists a constant $ C_s > 0 $ such that for any $ v \in H^s (\mathbb{T}) $, \label{p:s}
  \begin{align}
    \norm{0, 2}{(I - \mathcal{I}_{h}) v} & \leq C_s h^s \snorm{s, 2}{v}. \label{interpolation_error_s}
  \end{align}
  \item For an integer $ q > s $, there exists a constant $ C_q > 0 $ such that for any $ v \in H^q (\mathbb{T}) $, \label{p:q} 
  \begin{align} 
    \snorm{s, 2}{(I - \mathcal{I}_{h}) v} & \leq C_q h^{q - s} \snorm{q, 2}{v}. 
  \end{align} 
\end{enumerate} 

\begin{lemma} \label{lem:s} 
  Let $ v \in H^s (\mathbb{T}) $ and $ f \in W^{s, \infty} (\mathbb{T}) $ for an integer $ s \geq 2 $. We define $ \Delta t_1 = \min\{(3 \snorm{1, \infty}{f} \snorm{1, \infty}{v})^{-1}, 1\} $. If $ \Delta t < \Delta t_1 $, then there exists a unique $ w \in H^s (\mathbb{T}) $ such that 
  \begin{align} 
    w = v \circ {X_{\Delta t}^1}[f \circ w]. \label{a1} 
  \end{align} 
  Furthermore, $ {X_{\Delta t}^1}[f \circ w] $ is a bijection. 
\end{lemma} 

\begin{align} 
  \tau_h^n & = \frac{1}{\Delta t} \left\{u^n - (\mathcal{S}_{\Delta t}^{\mathrm{V}} u^{n - 1}) \circ {X_{\Delta t}^1} [f \circ u^n]\right\} \\ 
  & = \frac{1}{\Delta t} \left\{u^n - \frac{1}{2} ( u^{n - 1} \circ {X_{\Delta t}^1} [f \circ u^n - \tilde{\delta}] + u^{n - 1} \circ {X_{\Delta t}^1} [f \circ u^n + \tilde{\delta}]) \right\}, \label{def_tau} 
\end{align} 

\begin{align} 
  \norm{0, 2}{\tau_h^n} & \leq C_{\tau,0} \Delta t. \label{tau_0} 
\end{align} 

\begin{align} 
  \snorm{s, 2}{\tau_h^n} & \leq C_{\tau,s} \Delta t. \label{tau_s} 
\end{align} 

\begin{align} 
  \norm{s, 2, \ast}{\mathcal{S}_{\Delta t}^{\mathrm{V}} v} \leq \frac{1}{2} \left(\norm{s, 2, \ast}{v \circ {X_{\Delta t}^1}[- \tilde{\delta}]} + \norm{s, 2, \ast}{v \circ {X_{\Delta t}^1}[+ \tilde{\delta}]}\right) = \norm{s, 2, \ast}{v} \label{b7} 
\end{align} 

\begin{align} 
  \norm{s + 1, \infty}{\mathcal{S}_{\Delta t}^{\mathrm{V}} v} & \leq \frac{1}{2} \left(\norm{s + 1, \infty}{v \circ {X_{\Delta t}^1}[- \tilde{\delta}]} + \norm{s + 1, \infty}{v \circ {X_{\Delta t}^1}[+ \tilde{\delta}]}\right) \\ & = \norm{s + 1, \infty}{v}. \label{b7-2} 
\end{align} 

\begin{lemma} \label{lem:2} 
  Suppose that the interpolation operator $ \mathcal{I}_{h} $ possesses the properties (P\ref{p:o}) and (P\ref{p:s}) for some integer $ s \leq 2 $. Furthermore, suppose that there exists a positive constant $ C_{\Delta 1} $ such that $ h^{2s} / \Delta t \leq C_{\Delta 1} $. Let $ C_I = (C_s)^2 C_{\Delta 1} $, where $ C_s $ is the constant introduced in (P\ref{p:s}). Then we have for any $ v \in H^s (\mathbb{T}) $ and $ (h, \Delta t) \in (0,1)^2 $
  \begin{align} 
    \norm{s, 2, *}{\mathcal{I}_{h} v} & \leq \left(1 + C_I \Delta t \right) \norm{s, 2, *}{v}. 
  \end{align} 
\end{lemma} 

\begin{lemma} \label{lem:s2} 
  Let $ v \in H^s (\mathbb{T}) $ and $ f \in W^{s, \infty} (\mathbb{T}) $ for an integer $ s \geq 2 $. Let $ \Delta t_1 $ be as defined in Lemma~\ref{lem:s}. Then there exists a positive constant $ C_{\mathrm{P}2} $ which depends polynomially on $ \norm{s, \infty}{f} $ and $ \norm{s,2,\ast}{v} $ such that if $ \Delta t < \Delta t_1 $, 
  \begin{align} 
    \norm{s,2,\ast}{\mathcal{S}_{\Delta t}^{\mathrm{A}} (v)} \leq \left(1 + C_{\mathrm{P}2} \Delta t \right) \norm{s,2,\ast} {v}. \label{a2} 
  \end{align} 
\end{lemma} 

\begin{lemma}\label{lem:a1} 
  Let $ v \in H^s (\mathbb{T}) $ and $ f \in W^{s, \infty} (\mathbb{T}) $ for some integer $ s \geq 2 $. Then there exists a positive constant $ C_{\mathrm{P}3} $ which depends polynomially on $ \norm{s, 2, \ast}{v} $ such that 
  \begin{align} 
    \norm{s, 2, \ast}{f \circ v} \leq C_{\mathrm{P}3} \norm{s, \infty}{f}. \label{ap1} 
  \end{align} 
  
  Moreover, let $ v_1 $, $v_2 \in H^s (\mathbb{T}) $ and $ f \in W^{s + 1, \infty} (\mathbb{T}) $ for some integer $ s \geq 2 $. Then there exists a positive constant $ C_{\mathrm{P}4} $ which depends polynomially on $ \norm{s - 1, 2}{\partial_x v_1} $ and $ \norm{s - 1, 2} {\partial_x v_2} $ such that 
  \begin{align} 
    \norm{s, 2, \ast}{f \circ v_1 - f \circ v_2} \leq C_{\mathrm{P}4} \norm{s + 1, \infty}{f} \norm{s, 2, \ast}{v_1 - v_2}. \label{ap2} 
  \end{align} 
\end{lemma}

\begin{lemma}\label{lem:c} 
  Let $ v, w \in H^s (\mathbb{T}) $ and $ f \in W^{s,\infty} (\mathbb{T}) $ for $ s \geq 2 $. 
  Suppose that $ w $ satisfies $ \snorm{1, \infty}{w} \leq 1 / (2 \Delta t \snorm{1, \infty}{f}) $. 
  Then there exists a positive constant $ C_{\mathrm{P}5} $ which depends polynomially on $ \norm{s, \infty}{f} $ and $ \norm{s, 2}{w} $ such that 
  \begin{align} 
    \norm{s, 2, \ast}{v \circ {X_{\Delta t}^1}[f \circ w]} \leq (1 + C_{\mathrm{P}5} \Delta t) \norm{s, 2, \ast}{v}. \label{c2} 
  \end{align} 
  Furthermore, suppose there exists a positive constant $ C_{\Delta 1} $ such that $ h^{2s} / \Delta t $. Then there exists a positive constant $ C_{\mathrm{P}6} $ which depends polynomially on $ \norm{s, \infty}{f} $ and $ \norm{s, 2}{w} $ such that 
  \begin{align} 
    \norm{s, 2, \Delta}{v \circ {X_{\Delta t}^1}[f \circ w]} \leq (1 + C_{\mathrm{P}6} \Delta t) \norm{s, 2, \Delta}{v}. \label{c2-2} 
  \end{align} 
\end{lemma} 

  \begin{align}
    & \norm{s, 2, \ast}{v \circ {X_{\Delta t}^1} [f \circ w_1] - v \circ {X_{\Delta t}^1} [f \circ w_2]} \\ & \leq C_{\mathrm{P}7} \Delta t \norm{s + 1, \infty}{v} \norm{s, 2, \ast}{w_1 - w_2}. \label{d4} 
  \end{align} 

\begin{theorem} \label{thm:1} 
  Suppose that the interpolation operator $ \mathcal{I}_{h} $ satisfies (P\ref{p:o}), (P\ref{p:s}) and (P\ref{p:q}) for positive integers $ s $ and $ q $. Furthermore, assume that there exist positive constants $ C_{\Delta 2} $ and $ \varepsilon $ such that 
  \begin{align} 
    h^{2s} / \Delta t \leq C_{\Delta 2} \quad \text{and} \quad h^{2(q - s)} / \Delta t^{1 + \varepsilon} \leq C_{\Delta 2}. \label{hyp:delta} 
  \end{align} 
  Then there exist $ (h_0, \Delta t_0) \in (0, 1)^2 $ and $ C_e > 0 $ which depends on $ \norm{W^{1,\infty}(H^{s + 2})}{u} $, $ \norm{L^{\infty} (H^{s + 4})}{u} $ and $ \norm{L^{\infty}(H^{q})}{u} $ such that, if $ (h, \Delta t) \in (0, h_0) \times (0, \Delta t_0) $, the following estimate holds: 
  \begin{align} 
    \norm{s, 2, \ast}{u^n - u_h^n} \leq C_e \left(\Delta t + \frac{h^{q - s}}{\Delta t^{1/2}} + \frac{h^q}{\Delta t}\right) \label{b30} 
  \end{align} 
  for all $ n \in \{0, \ldots, N_t\} $. 
\end{theorem} 

\begin{proof} 
  In this proof, we do not explicitly indicate the dependence of the generic constant $C_{\mathrm{P}}$ on $f$. 

  Let $ e_h^n = u^n - u_h^n $. For $ n = 0 $, we have 
  \begin{align}
    \norm{s, 2, \ast}{e_h^0} = \left(\norm{0, 2}{(I - \mathcal{I}_{h}) u^0}^2 + \snorm{s, 2}{(I - \mathcal{I}_{h}) u^0}^2\right)^{1/2}. 
  \end{align}
  By (P\ref{p:s}) and (P\ref{p:q}), there exists a positive constant $ C_0 $ depending on $ \snorm{q, 2}{u} $ such that 
  \begin{align} 
    \norm{s, 2, \ast}{e_h^0} \leq C_0 h^{q - s}. 
    \label{b0}
  \end{align}
  For $ n \geq 2 $, by the second equality of \eqref{SL-V}, we have 
  \begin{align} 
    e_h^n 
    & = u^n - \mathcal{I}_{h} (\mathcal{S}_{\Delta t}^{\mathrm{A}} (\mathcal{S}_{\Delta t}^{\mathrm{V}} u_h^{n - 1})) \\
    & = \mathcal{I}_{h} (u^n - \mathcal{S}_{\Delta t}^{\mathrm{A}} (\mathcal{S}_{\Delta t}^{\mathrm{V}} u_h^{n - 1})) + (I - \mathcal{I}_{h}) u^n. 
    \label{b6}
  \end{align} 
  By \eqref{def_tau}, we can decompose the first term in the right-hand side of \eqref{b6} as 
  \begin{align}
    u^n - \mathcal{S}_{\Delta t}^{\mathrm{A}} (\mathcal{S}_{\Delta t}^{\mathrm{V}} u_h^{n - 1}) = \eta_h^n + \theta_h^n + \Delta t \, \tau_h^n, \label{b9} 
  \end{align} 
  where 
  \begin{gather}
    \begin{aligned}
      \eta_h^n & = (\mathcal{S}_{\Delta t}^{\mathrm{V}} u^{n - 1}) \circ {X_{\Delta t}^1} [f \circ \mathcal{S}_{\Delta t}^{\mathrm{A}} (\mathcal{S}_{\Delta t}^{\mathrm{V}} u_h^{n - 1})] - \mathcal{S}_{\Delta t}^{\mathrm{A}} (\mathcal{S}_{\Delta t}^{\mathrm{V}} u_h^{n - 1}), 
    \end{aligned} \\ 
    \begin{aligned} 
      \theta_h^n & = (\mathcal{S}_{\Delta t}^{\mathrm{V}} u^{n - 1}) \circ {X_{\Delta t}^1} [f \circ u^n] - (\mathcal{S}_{\Delta t}^{\mathrm{V}} u^{n - 1}) \circ {X_{\Delta t}^1} [f \circ \mathcal{S}_{\Delta t}^{\mathrm{A}} (\mathcal{S}_{\Delta t}^{\mathrm{V}} u_h^{n - 1})], \label{b9-3} 
    \end{aligned} 
  \end{gather} 
  and $ \tau_h^n $ is defined by \eqref{def_tau}. We can rewrite $ \eta_h^n $ as 
  \begin{align}
    \eta_h^n & = (\mathcal{S}_{\Delta t}^{\mathrm{V}} u^{n - 1}) \circ {X_{\Delta t}^1} [f \circ \mathcal{S}_{\Delta t}^{\mathrm{A}} (\mathcal{S}_{\Delta t}^{\mathrm{V}} u_h^{n - 1})] \\ & \quad - (\mathcal{S}_{\Delta t}^{\mathrm{V}} u_h^{n - 1}) \circ {X_{\Delta t}^1} [f \circ \mathcal{S}_{\Delta t}^{\mathrm{A}} (\mathcal{S}_{\Delta t}^{\mathrm{V}} u_h^{n - 1})] \\
    & = (\mathcal{S}_{\Delta t}^{\mathrm{V}} e_h^{n - 1}) \circ {X_{\Delta t}^1} [f \circ \mathcal{S}_{\Delta t}^{\mathrm{A}} (\mathcal{S}_{\Delta t}^{\mathrm{V}} u_h^{n - 1})]. \label{b9-2} 
  \end{align}
  Let $ \rho_h^n = (I - \mathcal{I}_{h}) u^n $. Then \eqref{b6} becomes 
  \begin{align}
    e_h^n = \mathcal{I}_{h} (\eta_h^n + \theta_h^n + \Delta t \tau_h^n) + \rho_h^n. \label{b6-1}
  \end{align}
  By \eqref{p:o}, we have 
  \begin{align}
    \snorm{s, 2}{e_h^n}^2 & = \snorm{s, 2}{\mathcal{I}_{h} (\eta_h^n + \theta_h^n + \Delta t \tau_h^n)}^2 + \snorm{s, 2}{\rho_h^n}^2. \label{b6-2} 
  \end{align}
  By an elementary inequality
  \begin{align} 
    (a_1 + a_2 + a_3)^2 \leq (1 + \Delta t) a_1^2 + \frac{4}{\Delta t} (a_2^2 + a_3^2) 
  \end{align} 
   for $ a_1, a_2, a_3 \in \mathbb{R} $ and $ \Delta t \in (0, 1) $, \eqref{b6-2} becomes 
  \begin{align} 
    \snorm{s, 2}{e_h^n}^2 & = \left(1 + \Delta t \right) \snorm{s, 2}{\mathcal{I}_{h} \eta_h^n}^2 + \frac{4}{\Delta t} \left(\snorm{s, 2}{\mathcal{I}_{h} \theta_h^n}^2 + \snorm{s, 2}{\Delta t \mathcal{I}_{h} \tau_h^n}^2\right) + \snorm{s, 2}{\rho_h^n}^2. \label{b12-1}
  \end{align} 
  By \eqref{b6-1} and the inequality 
  \begin{align}
    (a_1 + a_2 + a_3 + a_4)^2 \leq (1 + \Delta t) a_1^2 + \frac{6}{\Delta t} (a_2^2 + a_3^2 + a_4^2) 
  \end{align} 
  for $ a_1, a_2, a_3, a_4 \in \mathbb{R} $ and $ \Delta t \in (0, 1) $, we have 
  \begin{align} 
    \norm{0, 2}{e_h^n}^2 & \leq (1 + \Delta t) \norm{0, 2}{\mathcal{I}_{h} \eta_h^n}^2 \\ & \quad + \frac{6}{\Delta t} \left(\norm{0, 2}{\mathcal{I}_{h} \theta_h^n}^2 + \norm{0, 2}{\Delta t \mathcal{I}_{h} \tau_h^n}^2 + \norm{0, 2}{\mathcal{I}_{h} \rho_h^n}^2\right). \label{b12-2} 
  \end{align} 
  Combining \eqref{b12-1} and \eqref{b12-2}, we obtain 
  \begin{align} 
    \norm{s, 2, \ast}{e_h^n}^2
    & \leq (1 + \Delta t) \norm{s, 2, \ast}{\mathcal{I}_{h} \eta_h^n}^2 \\ & \quad + \frac{6}{\Delta t} \left(\norm{s, 2, \ast}{\mathcal{I}_{h} \theta_h^n}^2 + \norm{s, 2, \ast}{\Delta t \mathcal{I}_{h} \tau_h^n}^2 + \norm{0, 2}{\rho_h^n}^2\right)+ \snorm{s, 2}{\rho_h^n}^2. 
  \end{align} 
  By Lemma~\ref{lem:2}, we have 
  \begin{align}
    \norm{s, 2, \ast}{e_h^n}^2 & \leq (1 + C \Delta t) \norm{s, 2, \ast}{\eta_h^n}^2 + \frac{6}{\Delta t} \norm{s, 2, \ast}{\theta_h^n}^2 \\
    & \quad + 6 \Delta t \norm{s, 2, \ast}{\tau_h^n}^2 + \frac{6}{\Delta t} \norm{0, 2}{\rho_h^n}^2 + \snorm{s, 2}{\rho_h^n}^2. \label{b12}
  \end{align} 

  We estimate $ \norm{s, 2, \ast}{\eta_h^n}^2 $. By Lemma~\ref{lem:s2} and \eqref{b7}, we have 
  \begin{align} 
    \norm{2, s, \ast}{\mathcal{S}_{\Delta t}^{\mathrm{A}} (\mathcal{S}_{\Delta t}^{\mathrm{V}} u_h^{n - 1})} & \leq \left(1 + C_{\mathrm{P}} (\norm{s, 2, \ast}{\mathcal{S}_{\Delta t}^{\mathrm{V}} u_h^{n - 1}}) \Delta t\right) \norm{s, 2, \ast}{\mathcal{S}_{\Delta t}^{\mathrm{V}} u_h^{n - 1}} \\ 
    & \leq C_{\mathrm{P}} (\norm{s, 2, \ast}{u_h^{n - 1}}), \label{b8} 
  \end{align} 
  where we used $ \Delta t \leq 1 $. By Lemmas~\ref{lem:c} and \ref{lem:a1}, we have % chktex 2
  \begin{align}
    \norm{s, 2, \ast}{\eta_h^n} 
    & = \norm{s, 2, \ast}{(\mathcal{S}_{\Delta t}^{\mathrm{V}} e_h^{n - 1}) \circ {X_{\Delta t}^1} [f \circ \mathcal{S}_{\Delta t}^{\mathrm{A}} (\mathcal{S}_{\Delta t}^{\mathrm{V}} u_h^{n - 1})]} \\
    & \leq \left(1 + C_{\mathrm{P}} (\norm{s, 2, \ast}{f \circ (\mathcal{S}_{\Delta t}^{\mathrm{A}} \mathcal{S}_{\Delta t}^{\mathrm{V}} u_h^{n - 1})})  \Delta t\right) \norm{s, 2, \ast}{\mathcal{S}_{\Delta t}^{\mathrm{V}} e_h^{n - 1}}. 
  \end{align} 
  Using \eqref{b7} and \eqref{b8}, we have 
  \begin{align} 
    \norm{s, 2, \ast}{\eta_h^n} & \leq \left(1 + C_{\mathrm{P}} (\norm{s, 2, \ast}{u_h^{n - 1}})  \Delta t\right) \norm{s, 2, \ast}{e_h^{n - 1}}. \label{b2} 
  \end{align} 

  Second, we estimate $ \norm{s, 2, \ast}{\tau_h^n} $. By \eqref{tau_0} and \eqref{tau_s}, we have 
  \begin{align} 
    \norm{s, 2, \ast}{\tau_h^n} \leq C_{\mathrm{P}} (\norm{W^{2, \infty}(H^{s})}{u}, \norm{L^{\infty}(H^{s+4})}{u}) \Delta t. \label{b28}
  \end{align} 

  Third, we estimate $ \frac{6}{\Delta t} \norm{0,2}{\rho_h^n}^2 + \snorm{s,2}{\rho_h^n}^2 $. 
  By (P\ref{p:s}) and (P\ref{p:q}), we have 
  \begin{align}
    \frac{6}{\Delta t} \norm{0, 2}{\rho_h^n}^2 + \snorm{s, 2}{\rho_h^n}^2 \leq C\left(\frac{h^{2q}}{\Delta t} + h^{2 (q - s)}\right)\snorm{q, 2}{u}^2. \label{b31}
  \end{align}

  Now we estimate $ \norm{s, 2, \ast}{\theta_h^n} $. By \eqref{b9-3} and \eqref{d4}, we have 
  \begin{align} 
    \norm{s, 2, \ast}{\theta_h^n} 
    & \leq C_{\mathrm{P}} (\norm{s + 1, \infty}{\mathcal{S}_{\Delta t}^{\mathrm{V}} u^{n - 1}}, \norm{s, 2}{u^n}, \norm{s, 2}{\mathcal{S}_{\Delta t}^{\mathrm{A}} (\mathcal{S}_{\Delta t}^{\mathrm{V}} u_h^{n - 1})}) \\ & \qquad \cdot \Delta t
    \norm{s, 2, \ast}{u^n - \mathcal{S}_{\Delta t}^{\mathrm{A}} (\mathcal{S}_{\Delta t}^{\mathrm{V}} u_h^{n - 1})}. \label{b10-4} 
  \end{align} 
  By \eqref{b8} and \eqref{b7-2}, this becomes 
  \begin{align} 
    & \norm{s, 2, \ast}{\theta_h^n} \\ & \leq C_{\mathrm{P}} (\norm{L^\infty (W^{s + 1, \infty})}{u}, \norm{s, 2}{u_h^{n - 1}}) \Delta t \norm{s, 2, \ast}{u^n - \mathcal{S}_{\Delta t}^{\mathrm{A}} (\mathcal{S}_{\Delta t}^{\mathrm{V}} u_h^{n - 1})}.  \label{b10} 
  \end{align} 
  % By Lemma~\ref{lem:a1} and \eqref{b8}, we have 
  % \begin{align} 
  %   & \norm{s, 2, \ast}{f \circ u^n - f\circ \mathcal{S}_{\Delta t}^{\mathrm{A}} (\mathcal{S}_{\Delta t}^{\mathrm{V}} u_h^{n - 1})} \\
  %   & \leq C_{\mathrm{P}} (\norm{s, 2, \ast}{u^n}, \norm{s, 2, \ast}{\mathcal{S}_{\Delta t}^{\mathrm{A}} (\mathcal{S}_{\Delta t}^{\mathrm{V}} u_h^{n - 1})}) \norm{s + 1, \infty}{f} \norm{s, 2, \ast}{u^n - \mathcal{S}_{\Delta t}^{\mathrm{A}} (\mathcal{S}_{\Delta t}^{\mathrm{V}} u_h^{n - 1})} \\ 
  %   & \leq C_{\mathrm{P}} (\norm{s, 2, \ast}{u^n}, \norm{s, 2, \ast}{u_h^{n - 1}}) \norm{s, \infty}{f} \norm{s, 2, \ast}{u^n - \mathcal{S}_{\Delta t}^{\mathrm{A}} (\mathcal{S}_{\Delta t}^{\mathrm{V}} u_h^{n - 1})}. \label{b19}
  % \end{align}
  % By \eqref{b9}, \eqref{b10} becomes 
  % \begin{align} 
  %   & \norm{s, 2, \ast}{\theta_h^n} \\
  %   & \leq C_{\mathrm{P}} (\norm{L^\infty (W^{s + 1, \infty})}{u}, \norm{s, 2}{u_h^{n - 1}}) \Delta t \left(\norm{s, 2, \ast}{\rho_h^n} + \norm{s, 2, \ast}{\theta_h^n} + \Delta t \norm{s, 2, \ast}{\tau_h^n}\right). \label{b10-3}
  % \end{align} 
  % By \eqref{b9}, \eqref{b19} becomes 
  % \begin{align} 
  %   & \norm{s, 2, \ast}{f \circ u^n - f\circ \mathcal{S}_{\Delta t}^{\mathrm{A}} (\mathcal{S}_{\Delta t}^{\mathrm{V}} u_h^{n - 1})} \\
  %   & \leq C_{\mathrm{P}} (\norm{s, 2, \ast}{u^n}, \norm{s, 2, \ast}{u_h^{n - 1}}) \norm{s, \infty}{f}\left(\norm{s, 2, \ast}{\rho_h^n} + \norm{s, 2, \ast}{\theta_h^n} + \Delta t \norm{s, 2, \ast}{\tau_h^n}\right). \label{b11} 
  % \end{align} 
  % We note that $ \norm{s, 2, \ast}{u^n} $ can be bounded by $ \norm{s, 2, \ast}{u^n} \leq C \norm{L^{\infty} (W^{s + 1, \infty})}{u} $. 
  % Combining \eqref{b10} and \eqref{b11}, we obtain that there exists a constant $ C_{\mathrm{P}1} $ which depends polynomially on $ \norm{L^{\infty} (0, T; W^{s + 1, \infty} (\mathbb{T}))}{u} $ and $ \norm{s, 2, \ast}{u_h^{n-1}} $ such that 
  By \eqref{b9} and \eqref{b10}, there exists a constant $ C_{\mathrm{P}1} $ which depends polynomially on $ \norm{L^\infty (W^{s + 1, \infty})}{u} $ and $ \norm{s, 2, \ast}{u_h^{n-1}} $ such that 
  \begin{align}
    \norm{s, 2, \ast}{\theta_h^n} \leq C_{\mathrm{P}1} \Delta t \left(\norm{s, 2, \ast}{\rho_h^n} + \norm{s, 2, \ast}{\theta_h^n} + \Delta t \norm{s, 2, \ast}{\tau_h^n}\right). 
    \label{b10-2}
  \end{align}
  Let $ \Delta t_2 = \min \{1/(2 C_{\mathrm{P}1}), 1\} $ and suppose $ \Delta t \leq \Delta t_2 $. 
  Then we have $ 1 / (1 - C_{\mathrm{P}1} \Delta t) \leq 1 + 2 C_{\mathrm{P}1} \Delta t $. 
  Therefore, \eqref{b10-2} becomes 
  \begin{align}
    \norm{s, 2, \ast}{\theta_h^n} \leq C_{\mathrm{P}1} \Delta t (1 + 2 C_{\mathrm{P}1} \Delta t) \left(\norm{s, 2, \ast}{\eta_h^n} + \Delta t \norm{s, 2, \ast}{\tau_h^n}\right). 
    \label{b13}
  \end{align}
  Applying \eqref{b2}, \eqref{b31} and \eqref{b13} to \eqref{b12}, we obtain 
  \begin{align} 
    \norm{s, 2, \ast}{e_h^n}^2 & \leq \left(1 + C_{\mathrm{P}} (\norm{L^{\infty} (W^{s + 1, \infty})}{u}, \norm{s, 2, \ast}{u_h^{n - 1}} ) \Delta t \right) \norm{s, 2, \ast}{e_h^{n - 1}}^2 \\ & \quad + C_{\mathrm{P}} (\norm{L^{\infty} (W^{s + 1, \infty})}{u}, \norm{s, 2, \ast}{u_h^{n - 1}}) \Delta t \norm{s, 2, \ast}{\tau_h^{n}}^2 \\ & \quad + C\left(\frac{h^{2q}}{\Delta t} + h^{2 (q - s)}\right)\snorm{q, 2}{u}^2. \label{b20} 
  \end{align} 
  By \eqref{hyp:delta}, we estimate the last term as 
  \begin{align} 
    \left(\frac{h^{2q}}{\Delta t} + h^{2 (q - s)}\right)\snorm{q, 2}{u}^2 \leq (C_{\Delta 2}^2 + C_{\Delta 2}) \Delta t^{1 + \varepsilon}\snorm{q, 2}{u}^2. \label{b27} 
  \end{align} 
  Applying \eqref{b28} and \eqref{b27} to \eqref{b20}, we obtain 
  \begin{align} 
    \norm{s, 2, \ast}{e_h^n}^2 & \leq \left(1 + C_{\mathrm{P}} (\norm{L^{\infty} (W^{s + 1, \infty})}{u}, \norm{s, 2, \ast}{u_h^{n - 1}} ) \Delta t \right) \norm{s, 2, \ast}{e_h^{n - 1}}^2 \\ 
    & \quad + C_{\mathrm{P}} (\norm{W^{2, \infty}(H^{s})}{u}, \norm{L^{\infty}(H^{s+4})}{u}, \norm{s, 2, \ast}{u_h^{n - 1}}) \Delta t^3 \\ & \quad + C \Delta t^{1 + \varepsilon}\snorm{q, 2}{u}^2. \label{b18}
  \end{align} 
  Thus, there exist positive constants $ \alpha_1 $ depending on $ \norm{L^\infty(W^{s+1, \infty})}{u} $ and $ \beta_1 $ depending on $ \norm{W^{2, \infty}(H^s)}{u} $, $ \norm{L^{\infty}(H^{s + 4})}{u} $ and $ \norm{L^{\infty}(H^q)}{u} $ such that, if $ \norm{s, 2, \ast}{u_h^{n - 1}} \leq 2 \sup_{t\in (0, T)} \norm{s, 2, \ast}{u (t)} $, then 
  \begin{align} 
    \norm{s, 2, \ast}{e_h^n}^2 & \leq \left(1 + \alpha_1 \Delta t \right) \norm{s, 2, \ast}{e_h^{n - 1}}^2 + \beta_1 \Delta t^{1 + \varepsilon}. \label{b15} 
  \end{align} 
  Let positive constants $ \Delta t_3 $ and $ h_2 $ satisfy 
  \begin{align}
    \beta_1 T \mathrm{e}^{\alpha_1 T} (\Delta t_3)^\varepsilon \leq \frac{1}{2} \sup_{t\in (0, T)} \norm{s, 2, \ast}{u (t)} \label{b16} 
  \end{align} 
  and 
  \begin{align} 
    C_0 \mathrm{e}^{\alpha_1 T} (h_0)^{2(q - s)} \leq \frac{1}{2} \sup_{t\in (0, T)} \norm{s, 2, \ast}{u (t)}, \label{b17} 
  \end{align} 
  where $ C_0 $ is the constant introduced in \eqref{b0}. 
  
  Let $ \Delta t_0 = \min\{\Delta t_1, \Delta t_2, \Delta t_3, 1\} $ and $ h_0 = \min \{h_2, 1\} $. In the following, suppose that $ (\Delta t, h) \in (0, \Delta t_0) \times (0, h_0) $. We will show 
  \begin{align}
    \norm{s, 2, \ast}{u_h^n} \leq 2 \sup_{t\in (0, T)} \norm{s, 2, \ast}{u (t)}
    \label{b14}
  \end{align}
  for $ n \in \{0, \ldots, N_t\} $ by induction. 

  By \eqref{b0} and \eqref{b17}, we can easily check \eqref{b14} for $ n = 0 $ as 
  \begin{align} 
    \norm{s, 2, \ast}{u_h^0} \leq \norm{s, 2, \ast}{u^0} + \norm{s, 2, \ast}{e_h^0} \leq \norm{s, 2, \ast}{u^0} + C_0 h^{2 (q - s)} \leq \frac{3}{2} \sup_{t\in (0, T)} \norm{s, 2, \ast}{u (t)}. 
  \end{align}
  Suppose that, for some $ k \in \{0, \ldots, N_t - 1\} $, \eqref{b14} holds for $ n = 0, \ldots, k $. 
  Then by \eqref{b15}, we have 
  \begin{align} 
    \norm{s, 2, \ast}{e_h^{n + 1}}^2 & \leq \left(1 + \alpha_1 \Delta t \right) \norm{s, 2, \ast}{e_h^n}^2 + \beta_1 \Delta t^{1 + \varepsilon} 
  \end{align} 
  for all $ n = 0, \ldots, k $. From this, it follows that 
  \begin{align}
    \norm{s, 2, \ast}{e_h^{k + 1}}^2 
    & \leq \left(1 + \alpha_1 \Delta t \right)^{k + 1} \norm{s, 2, \ast}{e_h^0}^2 
    + \sum_{n = 0}^{k} \left(1 + \alpha_1 \Delta t \right)^{n} \beta_1 \Delta t^{1 + \varepsilon}. 
    \label{b29}
  \end{align}
  Applying \eqref{b0} and the inequality 
  \begin{align}
    \left(1 + \alpha_1 \Delta t \right)^{n} \leq \left(1 + \alpha_1 \Delta t \right)^{T / \Delta t} \leq \mathrm{e}^{\alpha_1 T} \quad (0 \leq n \leq N_t)
  \end{align}
  to \eqref{b29}, we obtain 
  \begin{align} 
    \norm{s, 2, \ast}{e_h^{k + 1}}^2 
    & \leq C_0 \mathrm{e}^{\alpha_1 T} h^{2 (q - s)}
    + \beta_1 T \mathrm{e}^{\alpha_1 T} \Delta t^{\varepsilon}. 
  \end{align} 
  Using \eqref{b16} and \eqref{b17}, we have 
  \begin{align} 
    \norm{s, 2, \ast}{e_h^{k + 1}}^2 \leq \sup_{t\in (0, T)} \norm{s, 2, \ast}{u (t)}. 
  \end{align} 
  Thus, we obtain 
  \begin{align} 
    \norm{s, 2, \ast}{u_h^{k + 1}} \leq \norm{s, 2, \ast}{u^{k + 1}} + \norm{s, 2, \ast}{e_h^{k + 1}} \leq 2 \sup_{t\in (0, T)} \norm{s, 2, \ast}{u (t)}. 
  \end{align} 
  By induction, \eqref{b14} holds for all $ n = 0,\ldots, N_t $. 

  Using \eqref{b28} and \eqref{b14}, \eqref{b20} can be rewritten as 
  \begin{align} 
    \norm{s, 2, \ast}{e_h^n}^2 & \leq \left(1 + \alpha_1 \Delta t \right) \norm{s, 2, \ast}{e_h^{n - 1}}^2 + C \left(\Delta t^3 + \frac{h^{2q}}{\Delta t} + h^{2 (q - s)}\right), 
  \end{align} 
  where $ C $ depends on $ \norm{W^{2, \infty}(H^{s})}{u} $, $ \norm{L^{\infty}(H^{s+4})}{u} $ and $ \norm{L^{\infty}(H^{q})}{u} $. From this, we can deduce 
  \begin{align} 
    \norm{s, 2, \ast}{e_h^n}^2 & \leq \left(1 + \alpha_1 \Delta t \right)^{n + 1} \norm{s, 2, \ast}{e_h^0}^2 + \sum_{k = 0}^{n} C \left(\Delta t^3 + h^{2 (q - s)} + \frac{h^{2q}}{\Delta t}\right) \left(1 + \alpha_1 \Delta t \right)^k \\ 
    & \leq C_0 \mathrm{e}^{\alpha_1 T} h^{2 (q - s)} + C T \mathrm{e}^{\alpha_1 T} \left(\Delta t^2 + \frac{h^{2 (q - s)}}{\Delta t} + \frac{h^{2q}}{\Delta t^2}\right), 
  \end{align} 
  which implies \eqref{b30}. 
\end{proof} 

\end{document}
